\begin{abstract}
\thispagestyle{empty}

Наименование работы: выпускная квалификационная работа на тему «Разработка клиент‑серверной системы корпоративного обучения с поддержкой адаптивного тестирования».

Цель работы — разработать клиент‑серверную систему корпоративного обучения с поддержкой адаптивного тестирования для персонализации учебного процесса сотрудников и студентов. Способы достижения: анализ предметной области и аналогов, формирование требований и сценариев использования, проектирование архитектуры и модели данных, описание алгоритма адаптивного тестирования, подготовка рекомендаций по внедрению и безопасности.\cite{moodle_lms,mobile_learning_solutions_2024,educause_adaptive_2016,yojji_multitenant_lms_2025}

Практический результат — прототип LMS с мультиарендной (multi-tenant) архитектурой, адаптивным тестированием и удобным интерфейсом для мобильного и веб‑доступа. Новизна и ценность состоят в комплексной проработке мобильного UX, встроенной адаптивной модели подбора заданий и проектировании единого серверного решения для нескольких организаций.\cite{metodichka_vkr_ikt_2021}

Работа состоит из введения и двух глав (предметная область и технологии; проектирование системы), списка использованных источников и приложений. Во введении обоснованы актуальность, цель и задачи работы, объект и предмет исследования. В первой главе описаны предметная область, постановка задачи, анализ аналогов и выбор технологий. Во второй главе представлены сценарии использования, архитектура и модель данных (диаграммы компонентов, процессов и БД), описание двух алгоритмов (адаптивный подбор вопросов и изоляция данных по тенантам), а также макеты основных страниц (прототипы интерфейсов).

Библиографический список состоит из источников по мере их упоминания в тексте (печатные и электронные в соответствии с ГОСТ 7.05-2008). Результаты работы могут быть апробированы в рамках внедрения прототипа в учебный процесс или представлены в виде доклада на конференции.

\end{abstract}
